O operador linear \(A\) em um espaço vetorial de dimensão finita \(V\) é chamado de involução se \(A^{2}=I\), onde \(I\) é o operador identidade. Seja \(\dim V=n\).
\begin{itemize}
    \item[(a)] Prove que para toda involução \(A\) em \(V\), existe uma base de \(V\) consistindo de autovetores de \(A\).
    \item[(b)] Determine o número máximo de involuções distintas que comutam entre si em \(V\).
\end{itemize}
